% !TEX root = ../../../tree_notes.tex
% =====================================================================
%  11.4  Compressed and Persistent Tries        sec:c08-variants
%  Exports: def:c08-radix  def:c08-persistent  lst:c08-persist
%           prob:c08-persistxor  lst:c08-persistxor-soln  ex:c08-radixlab  (internal lab, no URL)
%           (a plain comparison table, no label; Aho-Corasick notebox preview)
% =====================================================================
\section{Compressed and Persistent Tries}\label{sec:c08-variants}

\begin{previouslybox}
\begin{enumerate}[nosep, leftmargin=1.6em]
  \item In \defref{c08-nodelayout} a fixed \code{child[26]} array wasted space on lone children. What
        shape does the \texttt{care}/\texttt{cares} spine have -- how many children does each of $c,a,r,e$
        have? (exactly one each -- a \emph{chain} ripe for collapsing)
  \item \secref{c08-binary} added a per-node $\mathit{cnt}$ so a \emph{live} set could answer ``is this
        branch non-empty?'' What breaks if a query must instead ask about a set as it was $k$ insertions
        ago, while we keep inserting? (a single mutable count forgets the past -- we need \emph{versions})
  \item A persistent/functional data structure shares untouched parts between versions and copies only
        what changes. For a root-to-leaf trie insert, how many nodes actually change? (the ones on that
        one path -- depth-many, everything else is shared)
  \item Which prior chapter treats persistence as a general technique you can lean on for the full
        functional-array machinery? (Chapter~21, \chref{c19})
\end{enumerate}
\end{previouslybox}

\begin{topicintro}
The first three sections used one plain trie and never worried about its size or its history. This section
engineers the structure for the two situations where the naive form fails. \textbf{Compression}: a chain of
single-child nodes -- the $c$-$a$-$r$-$e$ spine shared by \texttt{car}, \texttt{care}, \texttt{cares} --
stores one letter per node and one fat array per node for information that is really a \emph{substring};
collapsing each maximal single-child chain into one edge yields the \term{radix} (compressed) trie
(\defref{c08-radix}), the same idea that makes a big alphabet affordable and that lives inside production
hash tables. \textbf{Persistence}: when queries reference the set \emph{as it was} at some moment -- the
first $k$ insertions, or ``values that existed before index $i$'' -- a single mutable count is not enough;
we keep a \term{version root} per prefix of insertions, \textbf{path-copying} only the depth-many nodes
that change and sharing the rest (\defref{c08-persistent}).

The section's solved anchor, LeetCode~1707 \probref{c08-persistxor}, is the purest persistence payoff:
each query $[x,m]$ wants $\max_y (x\oplus y)$ over the sub-set $\{y \in \texttt{nums}: y \le m\}$. Sort the
numbers, build a \emph{persistent binary trie} with one root per count of inserted values, and a query
reduces to a binary search (how many values are $\le m$?) followed by the \secref{c08-binary} complement
walk \emph{on that version's root} -- an $\bigO{\log C}$ answer that neither a mutable trie nor a fresh
rebuild gives cheaply. We also preview Aho--Corasick in one paragraph and leave the full persistence
abstraction to Chapter~21 (\chref{c19}).
\textbf{Prerequisites.} \secref{c08-triebasics} (layout/memory), \secref{c08-binary} (the complement greedy
and the $\mathit{cnt}$ field); binary search; a little pointer-sharing / ``copy-on-write'' intuition. No
prior exposure to persistence is assumed -- this section builds one working instance.

\textbf{After this section you can}
\begin{itemize}[nosep]
  \item explain what a radix trie collapses and why it saves memory on sparse/large alphabets; \hfill\lv{4}
  \item build a persistent binary trie by path-copying and query a specific version root; \hfill\lv{5}
  \item solve ``max XOR with a value $\le m$'' by combining a sorted prefix with version roots; \hfill\lv{5}
\end{itemize}
\end{topicintro}

\begin{recognitionbox}[Pattern card -- a trie that must not forget, or must not bloat]
\reckey{either (compress) the trie is \emph{too big} -- huge alphabet, long single-child chains, or a
memory limit the $|\Sigma|\cdot\text{nodes}$ array cannot meet -- or (persist) queries reference the set
at a \emph{point in its history}: a prefix of insertions, ``versions'', or a value range you turn into a
prefix by sorting.}
\recuse{
  \item ``$k$-th / max XOR over the first $i$ insertions, or over $[l,r]$ of a sorted key order'' -- a
        persistent trie, root per version (\probref{c08-persistxor}, \probref{c08-summit} kin);
  \item ``max XOR with $y \le m$'' -- sort keys, persistent root indexed by $\operatorname{upper\_bound}$
        (\probref{c08-persistxor});
  \item ``memory blew up on a plain trie / very large alphabet'' -- a radix/compressed trie
        (\defref{c08-radix});
  \item ``dictionary with long shared prefixes, need fast LCP'' -- compression again.
}
\recnotuse{
  \item a \emph{live}, history-free set with add/erase/query -- the mutable $\mathit{cnt}$ trie of
        \secref{c08-binary} is enough (and is exactly the Summit \probref{c08-summit});
  \item you need to \emph{edit} many past versions, not just read them -- full persistence with edits is
        heavier and belongs to the general treatment in \chref{c19};
  \item \emph{any} order statistic with deletion over a value domain -- a Fenwick/segment tree on
        compressed values (\chref{c13}, \chref{c14}) may beat a trie.
}
One-line test: \emph{is the question either ``too many nodes?'' (compress) or ``which version of the set?''
(persist)? Each maps to one of the two variants below.}
\end{recognitionbox}

\subsection{Intuition}

Two different pressures, two engineering answers.

\textbf{Memory / compression.} Look at \figref{c08-wordstrie} from \secref{c08-triebasics} once more. The
path to \texttt{care} is $c\!\to\!a\!\to\!r\!\to\!e$, and every one of those nodes has \emph{exactly one}
child -- they are a chain, not a decision. A 26-pointer array at each of them is mostly null: pure waste.
A \term{radix trie} stores the whole shared substring ``car'' on \emph{one} edge, keeping nodes only where
a real \emph{branch} occurs (or a word ends). Fewer nodes, fewer arrays, and lookup now compares a
substring instead of one letter -- still $\bigO{L}$ total, just fewer allocations. This is precisely why
$\Sigma=\{0,1\}$ binary tries rarely need it (max branching $2$, chains short) but letter tries over big
alphabets do.

\textbf{History / persistence.} The Summit and LC~1707 both need to ask ``over the set as it stood after
$c$ insertions, what is the max XOR?'' Rebuilding a trie per query is $\bigO{n}$ each. A \emph{persistent}
trie never destroys the old version: inserting a value \textbf{path-copies} -- it clones the $\le\!B$ nodes
on its root-to-leaf path (bumping their $\mathit{cnt}$) and \emph{shares} every untouched subtree with the
previous root (\defref{c08-persistent}). So $n$ versions cost $\bigO{nB}$ nodes total, but each version
root is a complete, independently queryable trie. Version root $k$ holds the first $k$ values; the
complement greedy of \secref{c08-binary} runs unchanged on it, only reading counts that are frozen for
that version.

\subsection{Definitions and Notation}

\begin{definition}{Radix (compressed / Patricia) trie}{c08-radix}
A \term{radix trie} collapses every maximal chain of single-child nodes into a single edge whose label is
the whole substring along that chain. Nodes survive only where the trie branches (a node with $\ge 2$
children) or where a key terminates. Membership/prefix lookups still run in $\bigO{L}$ (a substring
comparison per edge costs only the length it covers); the payoff is node and edge count proportional to
the number of \emph{keys and branches}, not to $\sum|{\text{key}}|$, which matters enormously for a large
or sparse alphabet. Rebuilding the compression after insertions costs a split operation; a \term{Patricia
trie} is the bitwise ($\Sigma=\{0,1\}$) specialisation.
\end{definition}

\begin{definition}{Persistent trie by path-copying}{c08-persistent}
A \term{persistent trie} keeps a \emph{root pointer per version}. Inserting into version $i$ to obtain
version $i{+}1$ does not mutate any existing node: it \textbf{clones} each node on the depth-$B$ path
being walked (copy its two/thirty-two child pointers and increment the cloned node's $\mathit{cnt}$),
points the new clone's relevant child at the next clone, and \emph{reuses} the unchanged child pointers
straight from version $i$. Cost is $\bigO{B}$ new nodes per insert and $\bigO{B}$ time; all untouched
subtrees are shared, so $n$ versions use $\bigO{nB}$ nodes rather than $\bigO{n^2 B}$. A query selects a
version by its root and runs the ordinary walk -- the frozen $\mathit{cnt}$ fields make it see exactly the
keys of that version. This is the read-only (``partial'') persistence the XOR problems need; general
persistence with edits is Chapter~21 (\chref{c19}).
\end{definition}

\begin{insightbox}
Path-copying is the whole trick, and it is cheap \emph{because a trie insert changes one root-to-leaf
path}. The number of new nodes per version is the depth ($B=31$ for integers), and every other node is
shared -- persistence on a trie is nearly free in a way persistence on, say, a balanced BST is not, since
a rotation there disturbs $\bigO{\log n}$ nodes with subtrees that themselves change shape. The
$\mathit{cnt}$ field from \secref{c08-binary} is what makes a shared subtree still answer ``is this branch
non-empty in \emph{this} version?'' correctly -- the count is per-version on the cloned spine and inherited
by reference elsewhere.
\end{insightbox}

\subsection{The Algorithm: A Working Persistent XOR-Trie}

Build roots $r_0, r_1, \dots, r_n$ over a sorted key list; $r_0$ is empty and $r_k$ contains the first $k$
keys. \code{insert} clones the path and bumps counts; \code{query} descends a chosen version's root using
the complement greedy, treating a null or $\mathit{cnt}=0$ branch as unavailable.

\begin{lstlisting}[caption={Persistent binary trie for max XOR: path-copying insert and a version-rooted complement query}, label={lst:c08-persist}]
struct PNode { PNode* ch[2]{}; int cnt = 0; };            // cnt = keys (in THIS version) passing here
static const int MAXB = 30;                               // values <= 1e9 -> bits 30..0

// Clone node n (or a fresh empty node if n is null): the essence of path-copying.
PNode* clone(PNode* n) { PNode* c = new PNode(); if (n) { c->ch[0]=n->ch[0]; c->ch[1]=n->ch[1]; c->cnt=n->cnt; } return c; }

// Return the root of version (prev + value x). Only the depth-(MAXB+1) path is copied.
PNode* insert(PNode* prev, int x, int b) {
    PNode* cur = clone(prev);
    cur->cnt++;
    if (b >= 0) {
        int k = (x >> b) & 1;
        PNode* pchild = prev ? prev->ch[k] : nullptr;
        cur->ch[k] = insert(pchild, x, b - 1);            // recurse down one branch; other stays shared
    }
    return cur;
}

// max x XOR y over the keys present in this version-root; -1 if the version is empty.
int query(PNode* u, int x, int b) {
    if (!u || u->cnt == 0) return -1;                     // no value here in THIS version
    if (b < 0) return 0;
    int k = (x >> b) & 1, want = k ^ 1;
    PNode* alt = u->ch[want];
    if (alt && alt->cnt > 0) return (1 << b) | query(alt, x, b - 1);   // take complement branch
    return query(u->ch[k], x, b - 1);                      // else same-bit branch (guaranteed cnt>0)
}
\end{lstlisting}

\begin{complexitybox}
Time: $\bigO{B}$ per version created (one path copy) and $\bigO{B}$ per query -- so building $n$ versions
is $\bigO{nB}$ and answering $Q$ queries is $\bigO{Q B}$. Memory: $\bigO{nB}$ nodes ($B{+}1$ new clones
per insert, everything else shared) -- for $n\le 2\cdot 10^5$, $B=31$ that is about $6\cdot 10^{6}$ small
two-pointer nodes, the real budget to check. Verdict: this is the structure that makes ``max XOR over a
value range / a prefix / a version'' an $\bigO{\log C}$ answer instead of an $\bigO{n}$ rebuild.
\end{complexitybox}

\subsection{Worked Example (in-chapter lab, no URL)}

\begin{example}{Radix compression and subtree counting on a shared spine}{c08-radixlab}
\emph{Compression.} Insert $\{$car, care, cares, cat$\}$. A plain trie uses nodes
$c,a,r,e,s$ and $t$ -- six nodes, five of them ($c,a,r,e$ and the single-child links) on chains. The
radix trie stores the edge ``ca'' then branches into ``r'' (which continues ``es'' with terminals at
``r'' and ``re''... ) -- concretely it keeps a node only at the root, at the $r$/$t$ fork, and at each
word end, folding $a$,$e$,$s$ single-child steps into substring labels: same lookups, fewer allocations.

\emph{Counting.} With a per-node $\mathit{cnt}$ (\defref{c08-prefixcount}) over the \emph{plain} trie of
$\mathcal{W}$ (\secref{c08-triebasics}), the $c$-$a$ node has $\mathit{cnt}=4$, $c$-$a$-$r$ has $3$,
$c$-$a$-$t$ has $1$, $d$ has $1$ -- ``words with prefix car'' $=3$ is a single walk to $r$ and a read.
Radix or plain, the count semantics are identical; compression only changes how many nodes hold them.
\end{example}

\subsection{Variations and Twists}

\begin{itemize}[nosep]
  \item \textbf{Full vs partial persistence}: what we built is \emph{partial} (only the newest version is
        edited, older ones are read-only). Editing arbitrary old versions needs fat nodes or more machinery
        -- \chref{c19}. \hfill\lv{5}
  \item \textbf{Offline instead of persistent}: LC~1707 also yields to \emph{sorting queries by $m$} and
        inserting into one \emph{mutable} trie as $m$ grows -- no cloning, same $\bigO{(n+Q)\log C}$, less
        memory, but it only works because queries are known up front and independent. \hfill\lv{4}
  \item \textbf{Version = time of insertion}: to answer ``over indices $[l,r]$'' build roots by index order
        and use root $r$ minus root $l{-}1$ conceptually via the range count -- the persistent prefix-sum
        idea (\chref{c19}, \chref{c13}). \hfill\lv{5}
  \item \textbf{Radix rebuild cost}: inserting into a compressed trie may \emph{split} an edge; amortise it
        or fall back to a plain trie when inserts dominate. \hfill\lv{4}
\end{itemize}

\subsection{Common Pitfalls}

\begin{warningbox}[Trap 1 -- mutating a shared node in a persistent insert]
The one rule of path-copying: \emph{never} write into an existing version's node. If your \code{insert}
follows a child pointer and mutates it in place instead of cloning, every older version that shares that
subtree silently changes -- queries on root $k$ see values inserted after $k$. Clone \emph{before}
descending (\coderef{c08-persist}), bump the clone's count, and only then re-link.
\end{warningbox}
\begin{warningbox}[Trap 2 -- reading $\mathit{cnt}$ from the wrong version]
A version root's correctness rests on each cloned spine node carrying the count \emph{for that version}.
Reusing a stale $\mathit{cnt}$ (e.g.\ incrementing the parent clone but comparing a child count that
belongs to the newer version) makes ``is this branch non-empty?'' lie. Every clone copies then increments;
the shared (uncopied) children already had the right count for the values they hold.
\end{warningbox}
\begin{warningbox}[Trap 3 -- empty-range queries]
When no key satisfies the version's condition ($\mathit{cnt}=0$ at the root, or $m$ below every value),
\code{query} must return \emph{no answer} (LC~1707 wants $-1$). A greedy descent that ignores the empty
case walks off a null child and crashes or returns a spurious value -- guard on $\mathit{cnt}=0$ first.
\end{warningbox}

\subsection{Problem Lab}

The section's anchor is LeetCode~1707, \probref{c08-persistxor}: for each query $[x,m]$ return
$\max_{\texttt{nums}[j]\le m}(x\oplus \texttt{nums}[j])$, or $-1$ if none qualify. Sorting the numbers and
building a persistent binary trie -- one root per count of smallest values -- turns ``nums $\le m$'' into
``query version root $= \operatorname{upper\_bound}$ index'', solved by the \secref{c08-binary} complement
walk on that frozen root.

% ---------------------------------------------------------------------------
\solvedproblem{c08-persistxor}{Maximum XOR With an Element From Array}{LeetCode 1707}
  {https://leetcode.com/problems/maximum-xor-with-an-element-from-array/}{LC Hard · L5}
  {persistent binary trie, path-copying roots, sort + upper\_bound by m, version complement query, -1 empty}

\begin{problemstatement}[Maximum XOR With an Element From Array]
Given an array \texttt{nums} of non-negative integers and \texttt{queries[i] = [x\textsubscript{i}, m\textsubscript{i}]},
answer query $i$ with the \emph{maximum} value of $x_i \oplus \texttt{nums}[j]$ over all $j$ such that
$\texttt{nums}[j] \le m_i$; if no such $j$ exists the answer is $-1$. Constraints (official):
$1 \le |\texttt{nums}|, |\texttt{queries}| \le 2\cdot 10^{5}$ and $0 \le \texttt{nums}[j], x_i, m_i \le 10^{9}$.
\end{problemstatement}

\begin{attemptgate}[Attempt it yourself first -- \lv{5}, give yourself 30 minutes]
Each query restricts to $\{\text{nums} \le m\}$. Sort \texttt{nums}; then that subset is a \emph{prefix} of
the sorted order. A structure that answers max-XOR over ``the first $k$ inserted'' is exactly a persistent
trie root $k$.
\peek{What does ``nums[j] $\le m$'' become after you sort nums?}
\nudge{A count $k=\operatorname{upper\_bound}(\text{sorted},m)$, i.e.\ the first $k$ sorted values. Build
a persistent trie with root $r_k$ holding those $k$ values, then run the complement query on $r_k$.}
\fullhint{Sort \texttt{nums}. Create roots $r_0..r_n$ where $r_{t+1}=\text{insert}(r_t,\text{sorted}[t])$.
For query $[x,m]$: $k=\operatorname{upper\_bound}$; if $k=0$ answer $-1$; else answer \code{query($r_k$,x)}.
Every operation is $\bigO{31}$.}
\end{attemptgate}

\solstep{Understanding the Problem}
For each $(x,m)$ we maximise $x\oplus y$ over the \emph{value-filtered} subset $\{y\in\texttt{nums}: y\le m\}$.
The unconstrained max-XOR is \probref{c08-maxxor}'s greedy; the filter $\le m$ is what forces history. Sort
\texttt{nums} ascending; then $\{y\le m\}$ is a prefix of length $k=\operatorname{upper\_bound}(m)$. Querying
max-XOR over ``the set of the first $k$ sorted values'' is a version query -- persistent roots give it in
$\bigO{B}$ without a rebuild. Empty prefix ($k=0$) $\Rightarrow$ answer $-1$.

\solstep{Key Observations}
\begin{observation}
After sorting, the condition $y\le m$ becomes ``$y$ is among the first $k$ values'' for
$k=\operatorname{upper\_bound}$ -- turning a value filter into a \emph{version} index (\defref{c08-persistent}).
\end{observation}
\begin{observation}
A persistent binary trie root $r_k$ answers max-XOR over its $k$ keys with the same complement greedy of
\lemref{c08-greedybit}; each node's $\mathit{cnt}$ tells the walk which branches are non-empty \emph{in
version $k$}.
\end{observation}
\begin{observation}
Building all $n{+}1$ roots costs $\bigO{nB}$ nodes (path-copying shares everything off the inserted path),
so total time is $\bigO{(n+Q)B}$ with $B=30$ bits for $\le 10^9$.
\end{observation}

\solstep{Think It Through With Me}
\begin{thinkbox}
\deadend{For each query, scan nums for the max $y\le m$ XOR $x$.}
{$O(Q\cdot n)=4\cdot 10^{10}$ -- \re. Independent queries make this hopeless.}
\deadend{Rebuild a fresh trie for each query's filtered subset.}
{Correct but $O(Q\cdot n B)$ -- building is the bottleneck; you are re-doing shared work every time.}
\aha{Sort once, build $n{+}1$ persistent roots by successive path-copying insert, then each query is a
binary search plus a $\bigO{B}$ complement walk on the matching version root -- $O((n+Q)\log C)$.}
\end{thinkbox}

\solstep{Solution}
Sort \texttt{nums}. Set $r_0$ empty and build $r_{t+1}$ from $r_t$ by inserting sorted[$t$], using the
path-copying \code{insert} of \coderef{c08-persist}; this gives one root per prefix count. For a query
$(x,m)$, let $k$ be the count of values $\le m$ (a binary search). If $k=0$ emit $-1$; else emit the
\code{query} result on root $r_k$ with value $x$ -- the complement greedy reading $r_k$'s frozen counts.
That is \coderef{c08-persist}'s two routines plus a sort and a binary search.

\solstep{Proof of Correctness}
\begin{proofbox}[Version root $r_k$ yields the constrained max XOR]
By \defref{c08-persistent}, $r_k$ encodes exactly the first $k$ sorted values and shares all structure
outside the copied path, so it is a faithful trie of $\{y\le m\}$ (the sorted prefix of that size). By
\lemref{c08-greedybit} the complement walk on $r_k$ returns $\max_{y \text{ in } r_k} x\oplus y$, which is
$\max_{y\in\texttt{nums},\,y\le m} x\oplus y$. If $k=0$ no $y\le m$ exists and $-1$ is the required output
by the problem statement. Hence every query returns the specified value. \end{proofbox}

\solstep{Implementation}
\begin{lstlisting}[caption={LeetCode 1707 -- sort, build persistent roots, upper\_bound + version query}, label={lst:c08-persistxor-soln}]
class Solution {
    struct PNode { PNode* ch[2]{}; int cnt=0; } *null_=nullptr;
    static const int MAXB = 30;
    PNode* ins(PNode* prev, int x, int b){
        PNode* cur = new PNode();                       // clone-on-write: always make a fresh node
        if (prev){ cur->ch[0]=prev->ch[0]; cur->ch[1]=prev->ch[1]; cur->cnt=prev->cnt; }
        cur->cnt++;
        if (b>=0){ int k=(x>>b)&1; PNode* pc = prev? prev->ch[k]:nullptr; cur->ch[k]=ins(pc,x,b-1); }
        return cur;
    }
    int qry(PNode* u, int x, int b){
        if(!u || u->cnt==0) return -1;
        if(b<0) return 0;
        int k=(x>>b)&1, want=k^1;
        PNode* alt=u->ch[want];
        if(alt && alt->cnt>0) return (1<<b) | qry(alt,x,b-1);
        return qry(u->ch[k],x,b-1);
    }
public:
    vector<int> maximizeXor(vector<int>& nums, vector<vector<int>>& queries) {
        sort(nums.begin(), nums.end());
        vector<PNode*> roots(1, ins(null_, 0, -1));     // roots[0] = empty version
        roots[0]->cnt = 0;
        for (int t=0;t<(int)nums.size();++t) roots.push_back(ins(roots.back(), nums[t], MAXB));
        vector<int> ans;
        for (auto& q : queries){
            int x=q[0], m=q[1];
            int k = upper_bound(nums.begin(), nums.end(), m) - nums.begin();  // count of values <= m
            ans.push_back(k==0 ? -1 : qry(roots[k], x, MAXB));
        }
        return ans;
    }
};
\end{lstlisting}

\solstep{Trace}
\begin{tracebox}[Official sample: nums=(0,1,2,3,4), queries (3,1),(1,3),(5,6) -> 3,3,7]
sorted nums (0,1,2,3,4); roots $r_0..r_5$ built by prefix insertion.\\
query (3,1): $k$ = count(nums$\le$1) = 2 (values 0,1). query on $r_2$ with $x{=}3$: $3\oplus 0{=}3$, $3\oplus 1{=}2$ -> max 3.\\
query (1,3): $k$ = count(nums$\le$3) = 4 (0,1,2,3). query on $r_4$ with $x{=}1$: best $1\oplus 2{=}3$ -> 3.\\
query (5,6): $k$ = count(nums$\le$6) = 5 (all). query on $r_5$ with $x{=}5$: $5\oplus 2{=}7$ -> 7.\\
Answers = (3, 3, 7). \checkmark\ matches the official output exactly.
\end{tracebox}

\solstep{Complexity}
\begin{complexitybox}
Time: sort $\bigO{n\log n}$; building $n{+}1$ roots $\bigO{nB}$; each query a $\bigO{\log n}$ binary
search plus $\bigO{B}$ walk. With $B=30$ and $n,Q\le 2\cdot 10^5$ this is roughly
$6\cdot 10^{6}{+}$, comfortably within limits. Memory: $\bigO{nB}$ trie nodes (shared across versions) --
the number to sanity-check against the memory cap, each node two pointers and an $\mathit{int}$. Verdict:
\ac. (The offline-by-$m$ variant drops persistence and the clone cost entirely -- see Variations.)
\end{complexitybox}

\solstep{Pitfalls and Variants}
\begin{warningbox}
Live slips: (i) mutating a shared node during insertion, corrupting older versions (Trap 1) -- always
allocate a fresh node; (ii) answering with the wrong version root (using $r_n$ instead of $r_k$) -- then
the $y\le m$ filter is silently ignored -- \textbf{WA}; (iii) not special-casing $k=0$ -- should emit $-1$.
\emph{Variant (no persistence):} sort queries by $m$, walk a single \emph{mutable} trie inserting values
as $m$ advances, answer online -- identical $\bigO{(n+Q)B}$ and lower memory, valid \emph{only} because
queries are independent and known in advance.
\end{warningbox}

\begin{notebox}[Aho--Corasick, in one paragraph (previewed, not built)]
The plain trie answers ``is this a prefix?'' but not ``does \emph{any} dictionary word appear as a
\emph{substring} of this long text?''. \term{Aho--Corasick} adds a \emph{failure link} on every node -- a
longest-proper-suffix pointer, essentially KMP generalised from one pattern to a trie of many -- so a
single left-to-right scan of the text, following trie edges and falling back on failure links, reports
every dictionary hit in $\bigO{|\text{text}| + \text{matches}}$. It is a whole chapter's worth of BFS
construction and link semantics, explicitly \emph{out of scope} here: this chapter owns reading a key
symbol-by-symbol from the root, and Aho--Corasick is the multi-pattern substring extension of that idea.
\end{notebox}

\begin{center}
\renewcommand{\arraystretch}{1.2}
\setlength{\tabcolsep}{4pt}
\begin{tabular}{@{}p{0.24\linewidth} p{0.24\linewidth} p{0.46\linewidth}@{}}
\toprule
\textbf{Variant} & \textbf{Node cost / op} & \textbf{Use it when} \\
\midrule
Plain trie ($\Sigma$ letters) & $\bigO{|\Sigma|\cdot\text{nodes}}$ memory & small alphabet, history-free
  membership \& prefix queries (\secref{c08-triebasics}) \\
Binary (XOR) trie & $2$ children, depth $\bigO{B}$ & integer max/min XOR, live add/erase
  (\secref{c08-binary}, \probref{c08-summit}) \\
Radix / compressed & nodes $\propto$ keys $+$ branches & large/sparse alphabet, long shared prefixes
  (\defref{c08-radix}) \\
Persistent & $\bigO{B}$ new nodes per version & queries over a \emph{version / prefix / value range}
  (\defref{c08-persistent}, \probref{c08-persistxor}) \\
\bottomrule
\end{tabular}
\end{center}

\begin{insightbox}[Key takeaways]
\begin{enumerate}[nosep, leftmargin=1.6em]
  \item Two engineering axes on the same skeleton: \emph{memory} (compress single-child chains into a
        radix trie, \defref{c08-radix}) and \emph{time/history} (path-copying persistent roots,
        \defref{c08-persistent}).
  \item A persistent XOR-trie turns ``max XOR with $y\le m$ / within a range / after $k$ insertions'' into
        $\bigO{B}$ per version by keeping a frozen $\mathit{cnt}$ spine and sharing everything else.
  \item Never mutate a shared node: clone-then-write is the one rule that keeps old versions correct
        (Trap 1). General persistence, and Aho--Corasick's failure links, live in \chref{c19} and beyond.
\end{enumerate}
\end{insightbox}

\begin{hardeningbox}
\begin{itemize}[nosep, leftmargin=1.4em]
  \item \textbf{Cold-write \coderef{c08-persist}} and reproduce the official trace (3,3,7) of
        \probref{c08-persistxor} by hand on the five sorted roots.
  \item \textbf{Code the offline alternative}: sort queries by $m$, use one mutable trie from
        \coderef{c08-maxxor}; confirm identical answers without any cloning.
  \item \textbf{Read ahead}: the same path-copying gives persistent segment trees for range $k$-th --
        Chapter~21 (\chref{c19}); value-domain order statistics with deletes -- Chapters~13--14
        (\chref{c13}, \chref{c14}).
\end{itemize}
\end{hardeningbox}

\subsection{Practice Problems}

\begin{stuckbox}
\begin{itemize}[nosep, leftmargin=1.4em]
  \item Did you turn the value constraint ($y\le m$) into a version index via a sort + binary search?
  \item Does your insert always allocate a fresh node instead of editing a shared one?
  \item Did you return $-1$ when the filtered set is empty, rather than walking a null branch?
\end{itemize}
\end{stuckbox}

\begin{practiceset}
  \practice{Count Pairs With XOR in a Range}{LeetCode 1803}
    {https://leetcode.com/problems/count-pairs-with-xor-in-a-range/}{LC Hard · L5}
    {Count pairs $(i,j)$ with a low $\le$ nums[i]$\oplus$nums[j] $\le$ high. A binary trie answers
    ``how many stored $y$ give $x\oplus y < K$'' by descending MSB-first and, whenever $x$'s bit equals
    $K$'s bit, adding the whole \emph{same-bit} branch's $\mathit{cnt}$ before continuing on the path that
    keeps the prefix $\le K$; take $\text{count}(<\text{high}{+}1)-\text{count}(<\text{low})$. Insert
    elements progressively so each $x$ is queried only against earlier ones (no double count). It reuses
    the $\mathit{cnt}$ idea from \probref{c08-persistxor} on a \emph{mutable} trie -- persistence is not
    required because there is no value filter, just a live set.}
\end{practiceset}

\begin{recallbox}[10-second refresher -- 11.4 Compressed \& Persistent Tries]
Two ways the plain trie fails and their fixes. \textbf{Too big}: a radix/compressed trie folds maximal
single-child chains into one substring edge, keeping nodes only at branches/terminals (\defref{c08-radix}).
\textbf{Must not forget}: a persistent trie path-copies the $\bigO{B}$ nodes on an insert and shares the
rest, giving a root per version (\defref{c08-persistent}) -- the frozen $\mathit{cnt}$ spine makes a version
queryable. LC~1707 sorts, builds roots, and answers ``max XOR with $y\le m$'' by
$k=\operatorname{upper\_bound}$ then a complement walk on root $k$, or $-1$ (\probref{c08-persistxor}).
General persistence \chref{c19}; value-domain with deletes \chref{c13}, \chref{c14}; Aho--Corasick is the
failure-link preview, not taught.
\selftest{
  \item What does a radix trie collapse, and when does that save the most?
  \item State the one rule of path-copying and the bug you hit by breaking it.
  \item How does sorting turn ``$y\le m$'' into a persistent-trie version query?
  \item When is the \emph{offline} (mutable trie) solution to LC~1707 acceptable?
  \item What field lets a version root answer ``is this branch non-empty?'' for its frozen set?
}
\end{recallbox}
